% ============================================================
%  Computational indistinguishability and pseudorandom generators
% ============================================================

\section{Pseudorandom generators}

\subsection{Motivation: from a few random bits to many}

An extractor gives us a few truly random bits, but we often need many more: for example, a One-Time Pad key for a 1-gigabyte file. A \textbf{pseudorandom generator} (PRG) takes a short truly random string, the \textbf{seed}, and stretches it into a long string that \textbf{looks} random.

\textbf{Looks random to whom?} The output of a PRG \textbf{cannot} be truly uniform. If the seed has $n$ bits, the generator produces at most $2^n$ different strings out of the $2^{n+\ell}$ possible ones. An \textbf{unbounded} adversary can try all the seeds and check whether a given string is one of the possible outputs. So pseudorandomness makes sense \textbf{only against efficient adversaries}. This does not contradict Shannon's theorem, which holds against unbounded adversaries.

\subsection{Computational indistinguishability}

We need the computational counterpart of the statistical distance (\Cref{def:sd}). Since security is asymptotic in $n$, we talk about \textbf{probability ensembles}: sequences $\mathcal{X} = \{X_n\}_{n \in \N}$ of distributions, one for every value of the security parameter.

\begin{frameddefn}[label={def:cind}]{Computational indistinguishability}
  Two ensembles $\mathcal{X} = \{X_n\}_{n \in \N}$ and $\mathcal{Y} = \{Y_n\}_{n \in \N}$ are \textbf{computationally indistinguishable}, written $\mathcal{X} \cind \mathcal{Y}$, if for every PPT algorithm $D$ (the \textit{distinguisher}) there is a negligible function $\negl$ such that, for every $n$:
  \[
    \Big|\Pr[D(z) = 1 : z \leftarrow X_n] - \Pr[D(z) = 1 : z \leftarrow Y_n]\Big| \leq \negl(n)
  \]
\end{frameddefn}

It is the same ``quiz'' as always: $D$ receives a string and must say which distribution it comes from, and it must not do noticeably better than tossing a coin. The differences with the statistical distance are that the distinguisher must be \textbf{efficient}, and that the advantage must be \textbf{negligible} as $n$ grows.

Useful properties:
\begin{itemize}
  \item \textbf{Efficient processing cannot help} $\to$ if $\mathcal{X} \cind \mathcal{Y}$ and $f$ is computable in polynomial time, then $f(\mathcal{X}) \cind f(\mathcal{Y})$. Otherwise, a distinguisher for $f(\mathcal{X})$ and $f(\mathcal{Y})$ composed with $f$ would distinguish $\mathcal{X}$ from $\mathcal{Y}$.
  \item \textbf{Transitivity} $\to$ if $\mathcal{X} \cind \mathcal{Y}$ and $\mathcal{Y} \cind \mathcal{Z}$, then $\mathcal{X} \cind \mathcal{Z}$, because the sum of two negligible functions is negligible. This extends to a \textit{constant} number of steps; for a polynomial number of steps we need the hybrid argument of \Cref{thm:stretch}.
  \item \textbf{Statistical closeness implies computational indistinguishability} $\to$ if $\SD(X_n; Y_n)$ is negligible, then $\mathcal{X} \cind \mathcal{Y}$, since the statistical distance bounds the advantage of every distinguisher, efficient or not.
\end{itemize}

\subsection{The definition of PRG}

\begin{frameddefn}[label={def:prg}]{Pseudorandom generator}
  A function $G : \{0,1\}^n \to \{0,1\}^{n+\ell}$, with $\ell \geq 1$, is a \textbf{pseudorandom generator} with stretch $\ell$ if:
  \begin{enumerate}
    \item $G$ is computable in $\poly(n)$ time;
    \item $G(\Uni{n}) \cind \Uni{n+\ell}$, where $\Uni{k}$ is the uniform distribution over $\{0,1\}^k$.
  \end{enumerate}
\end{frameddefn}

Observations:
\begin{itemize}
  \item \textbf{$G$ is deterministic} $\to$ all the randomness of the output comes from the seed. The strength of the definition is that $n$ random bits become $n + \ell$ bits that no efficient algorithm can tell apart from uniform.
  \item \textbf{The seed must be uniform} (or at least computationally indistinguishable from uniform). In practice, a randomness extractor produces the seed, and the PRG expands it.
  \item \textbf{A first application.} With a PRG, Alice and Bob can share a short seed $s$ and use $G(s)$ as the pad of a One-Time Pad for a much longer message. The scheme is not perfectly secret (Shannon's theorem forbids it), but it is secure against efficient adversaries. We will make this precise in the next chapters.
\end{itemize}

\subsection{Stretching a PRG}

How much should a PRG stretch? Surprisingly, it is enough to build a PRG that adds \textbf{a single bit}: from it we can obtain a PRG that adds any polynomial number of bits.

\paragraph{The construction.} Let $G : \{0,1\}^n \to \{0,1\}^{n+1}$ be a PRG, and split its output as $G(s) = (s', b)$ with $s' \in \{0,1\}^n$ and $b \in \{0,1\}$. Apply $G$ to the seed, keep the last bit as output, and apply $G$ again to the remaining $n$ bits:
\[
  s_0 \xrightarrow{\ G\ } (s_1, b_1), \qquad s_1 \xrightarrow{\ G\ } (s_2, b_2), \qquad \ldots, \qquad s_{\ell-1} \xrightarrow{\ G\ } (s_\ell, b_\ell)
\]
and define
\[
  G^\ell(s_0) = b_1 \concat b_2 \concat \cdots \concat b_\ell \concat s_\ell \in \{0,1\}^{n+\ell}
\]

\begin{figure}[H]
  \centering
  \begin{tikzpicture}[g/.style={draw, rounded corners=2pt, minimum size=6mm}, x=1.15cm]
    \node (s0) at (0,0) {$s_0$};
    \node[g] (g1) at (1.2,0) {$G$};
    \node (s1) at (2.4,0) {$s_1$};
    \node[g] (g2) at (3.6,0) {$G$};
    \node (s2) at (4.8,0) {$s_2$};
    \node (dots) at (6,0) {$\cdots$};
    \node[g] (gl) at (7.2,0) {$G$};
    \node (sl) at (8.4,0) {$s_\ell$};
    \node (b1) at (1.2,-1.2) {$b_1$};
    \node (b2) at (3.6,-1.2) {$b_2$};
    \node (bl) at (7.2,-1.2) {$b_\ell$};
    \draw[msg] (s0) -- (g1);
    \draw[msg] (g1) -- (s1);
    \draw[msg] (s1) -- (g2);
    \draw[msg] (g2) -- (s2);
    \draw[msg] (s2) -- (dots);
    \draw[msg] (dots) -- (gl);
    \draw[msg] (gl) -- (sl);
    \draw[msg] (g1) -- (b1);
    \draw[msg] (g2) -- (b2);
    \draw[msg] (gl) -- (bl);
  \end{tikzpicture}
  \caption{The stretching construction: every application of $G$ outputs one bit and a new state.}
\end{figure}

\begin{framedthm}[label={thm:stretch}]{PRG stretching}
  If $G : \{0,1\}^n \to \{0,1\}^{n+1}$ is a PRG, then for every $\ell = \poly(n)$ the function $G^\ell : \{0,1\}^n \to \{0,1\}^{n+\ell}$ defined above is a PRG.
\end{framedthm}

The proof uses the \textbf{hybrid argument}, one of the most important techniques of the course. The idea is to build a chain of distributions that goes, one small step at a time, from the output of $G^\ell$ to a truly uniform string.

\paragraph{The hybrids.} For $i = 0, \ldots, \ell$, the hybrid $H_i$ is the distribution of the output of this procedure:
\begin{enumerate}
  \item choose the first $i$ output bits $b_1, \ldots, b_i$ uniformly at random, and a uniform state $s_i \sample \{0,1\}^n$;
  \item for $j = i+1, \ldots, \ell$, compute $(s_j, b_j) = G(s_{j-1})$;
  \item output $b_1 \concat \cdots \concat b_\ell \concat s_\ell$.
\end{enumerate}
In words, in $H_i$ the first $i$ bits are \textbf{truly random}, and from there on the string is produced by $G$.

\begin{table}[H]
  \centering
  \small
  \begin{tabular}{lccccc}
    \toprule
    \textbf{Hybrid} & $b_1$ & $b_2$ & $\cdots$ & $b_\ell$ & $s_\ell$ \\
    \midrule
    $H_0 = G^\ell(\Uni{n})$ & $G$ & $G$ & $\cdots$ & $G$ & $G$ \\
    $H_1$ & random & $G$ & $\cdots$ & $G$ & $G$ \\
    $H_2$ & random & random & $\cdots$ & $G$ & $G$ \\
    $\vdots$ & & & & & \\
    $H_\ell = \Uni{n+\ell}$ & random & random & $\cdots$ & random & random \\
    \bottomrule
  \end{tabular}
  \caption{The hybrids: ``$G$'' means produced by the chain of $G$, ``random'' means truly uniform. In $H_\ell$ the final state $s_\ell$ is the uniform state chosen in step 1.}
\end{table}

The two extremes are exactly the distributions we want to compare: $H_0 = G^\ell(\Uni{n})$ and $H_\ell = \Uni{n+\ell}$.

\begin{proof}
  $G^\ell$ is computable in polynomial time, since it runs $G$ exactly $\ell = \poly(n)$ times. We must show that $H_0 \cind H_\ell$.

  \textbf{Step 1: neighbouring hybrids differ in one application of $G$.} Compare $H_i$ and $H_{i+1}$. In both, $b_1, \ldots, b_i$ are random and $s_i$ is uniform. The only difference is the pair $(s_{i+1}, b_{i+1})$:
  \begin{itemize}
    \item in $H_i$ it is $G(s_i)$ with $s_i$ uniform, that is, a sample of $G(\Uni{n})$;
    \item in $H_{i+1}$ it is a uniform string of $n+1$ bits, that is, a sample of $\Uni{n+1}$.
  \end{itemize}
  Everything after that pair is computed in the same way, by applying $G$.

  \textbf{Step 2: a distinguisher for neighbours becomes a distinguisher for $G$.} Let $D$ be any PPT algorithm. For each $i$, define the algorithm $D_i$ that, on input $w = (s', b') \in \{0,1\}^{n+1}$:
  \begin{enumerate}
    \item chooses $b_1, \ldots, b_i$ uniformly at random;
    \item sets $s_{i+1} = s'$ and $b_{i+1} = b'$;
    \item computes $(s_j, b_j) = G(s_{j-1})$ for $j = i+2, \ldots, \ell$;
    \item returns $D(b_1 \concat \cdots \concat b_\ell \concat s_\ell)$.
  \end{enumerate}
  If $w \leftarrow G(\Uni{n})$, the string given to $D$ is distributed exactly as $H_i$; if $w \leftarrow \Uni{n+1}$, it is distributed exactly as $H_{i+1}$. Writing $p_i = \Pr[D(z) = 1 : z \leftarrow H_i]$:
  \[
    \Pr[D_i(w) = 1 : w \leftarrow G(\Uni{n})] - \Pr[D_i(w) = 1 : w \leftarrow \Uni{n+1}] = p_i - p_{i+1}
  \]

  \textbf{Step 3: from neighbours to the extremes.} Let $D^*$ be the algorithm that chooses $i \in \{0, \ldots, \ell-1\}$ uniformly at random and runs $D_i$. It is PPT because $\ell$ is polynomial. Its advantage against $G$ is the average of the $\ell$ differences, and the sum \textbf{telescopes}:
  \[
    \Pr[D^*(w) = 1 : w \leftarrow G(\Uni{n})] - \Pr[D^*(w) = 1 : w \leftarrow \Uni{n+1}]
    = \frac{1}{\ell}\sum_{i=0}^{\ell-1} (p_i - p_{i+1}) = \frac{p_0 - p_\ell}{\ell}
  \]
  Since $G$ is a PRG, the left side is at most $\negl(n)$ in absolute value. Therefore
  \[
    \big|\Pr[D(z) = 1 : z \leftarrow G^\ell(\Uni{n})] - \Pr[D(z) = 1 : z \leftarrow \Uni{n+\ell}]\big| = |p_0 - p_\ell| \leq \ell \cdot \negl(n)
  \]
  and $\ell \cdot \negl(n)$ is still negligible because $\ell$ is a polynomial. Since $D$ was arbitrary, $G^\ell(\Uni{n}) \cind \Uni{n+\ell}$.
\end{proof}

\begin{framedobs}{Why $\ell$ must be polynomial}
  The advantage between the two extremes is at most the \textbf{sum} of the advantages between neighbours. With $\ell$ hybrids this is at most $\ell \cdot \negl(n)$, which is negligible only if $\ell$ is polynomial: a polynomial times a negligible function is negligible, while an exponential times a negligible function may not be. This is why ``transitivity of $\cind$'' alone is not enough here: it would only cover a constant number of steps.
\end{framedobs}
